% !TeX root = ../main.tex

\chapter{結論與未來研究方向}

\section{研究結論}

本研究以知識圖譜注意力網路（KGAT）為核心，在 Food.com 食譜推薦資料集上進行效能評估與可解釋性分析。透過基準方法對比、深度比較、模組消融與 Fidelity 量化評估，本文對第一章所提出之三個研究問題提供了具一致趨勢的實證回答，並整理出若干與食譜知識圖譜特性相關的重要觀察。

需要指出的是，本文之結論主要建立於 Food.com 資料集、本文採用之評估協定，以及既定運算資源條件下的實驗結果。因此，下述發現較適合被理解為在本研究設定下所獲得之具體實證，而其跨資料集與跨設定的普適性，仍有待後續研究進一步驗證。

\subsection*{核心研究貢獻}

本研究的主要貢獻可歸納為三點：

\begin{enumerate}[leftmargin=2em]
  \item \textbf{深度效應之實證分析：}在食譜推薦領域觀察到 GNN 傳播深度與推薦效能之間存在顯著關聯，以從 $L{=}1$ 到 $L{=}3$ 的逐層消融實驗，量化了高階連接性的具體貢獻。
  \item \textbf{淺層知識圖譜資訊汙染現象之揭露與機制溯源：}發現並解釋了在 $L{=}1$ 淺層架構下，引入知識圖譜反而導致效能下降的「資訊汙染」現象。此反直覺發現填補了傳統認知中「引入知識必然有益」的盲點，並提出了具體的機制解釋——高連接度樞紐節點在淺層聚合中稀釋協同過濾訊號。
  \item \textbf{深層模型之端到端 Fidelity 驗證框架：}建立了一套從注意力路徑萃取到忠實度量化的完整可解釋性驗證管線，並在最佳效能的 $L{=}3$ 模型上對 103 條解釋路徑進行了初步 Fidelity 評估，揭露了 KG 語意路徑的充分性危機（$F^- = 0.3134$）——即人類語意上合理的解釋路徑，在數學上不足以支撐模型的預測信心。
\end{enumerate}

\subsection*{回應 RQ1：圖神經網路深度與知識圖譜整合如何影響推薦效能？}

實驗結果指出\textbf{訊息傳播深度是決定推薦效能的關鍵因素}。三層架構 ($L{=}3$) 的 KGAT 在所有指標上全面超越基準方法，其 HR@20 達到 0.9368，較最強的 Baseline (LightGCN, 0.6717) 提升了 39.5\%，NDCG@20 達到 0.5891，提升幅度更高達 52.7\%。深度每增加一層，HR@20 皆呈現顯著提升，充分支持在知識圖譜推薦場景中，二跳與三跳的高階連接性蘊含了極為關鍵的協同過濾訊號。

同時，消融實驗揭露了一個重要的反直覺發現：在淺層架構 ($L{=}1$) 下，移除知識圖譜後 HR@20 反而從 0.7910 提升至 0.8621（+9.0\%），NDCG@20 從 0.4440 提升至 0.5255（+18.4\%）。我們推測此現象的主因在於食譜知識圖譜中大量高連接度的樞紐實體節點——例如「chicken」或「breakfast」等常見食材與標籤——在 $L{=}1$ 的單跳聚合中，將高度泛化的特徵無差別地灌注至食譜嵌入，嚴重稀釋了純粹的協同過濾訊號。知識圖譜的引入需要搭配足夠的網路深度，使模型能透過多跳傳播逐步篩選有價值的語意關聯，同時以注意力機制抑制雜訊。

\subsection*{回應 RQ2：不同傳播深度下，模型的訓練行為與泛化能力有何差異？}

從附錄~\ref{app:training_log} 所呈現的 KGAT $L{=}3$ 模型三個完整 Epoch 訓練歷程來看，模型在訓練初期便展現了極快的收斂速度：第 1 個 Epoch 結束時，HR@20 已達到 0.8906，至第 3 個 Epoch 微幅上升至 0.9185，增幅僅 3.1\%。配合 NDCG@20 從 0.5406 上升至 0.5687（+5.2\%）來看，深層模型的核心效能在前三個 Epoch 內的增長空間已顯著收窄，可能預示著更長訓練週期下會面臨過擬合的風險。

此外，為進一步驗證模型在更長訓練週期的行為，本研究在偶數 Epoch 的模型檢查點上進行了額外的完整指標評估（詳見附錄~\ref{app:ablation_data}）。在此擴展評估窗口中，$L{=}3$ 的 HR@20 從 E2 的 0.9120 持續上升至 E10 的 0.9368，但增幅急劇趨緩（E2$\to$E10 僅增加 2.7\%，而 $L{=}1$ 在相同區間增加 6.2\%），再次印證深層模型的核心效能在訓練初期便已大致形成。

\subsection*{回應 RQ3：基於注意力機制的解釋路徑是否忠實地反映了模型的決策邏輯？}

依研究設計，抽取 100 位使用者、其中 60 位成功萃取解釋路徑進行初步忠實度評估，本研究驗證了注意力解釋路徑展現出\textbf{方向一致但程度差異化的忠實度}：

\begin{itemize}[leftmargin=2em]
  \item 93.3\% 的路徑 $F^+ > 0$，平均 $F^+ = 0.0711$，顯示注意力權重確實標記了對預測有貢獻的邊。
  \item 直接連接路徑的忠實度最高 (Avg $F^+ = 0.1231$, $F^- \approx 0$)，可被視為高可信度的解釋類型。
  \item 協同過濾路徑同樣具備良好忠實度 (Avg $F^+ = 0.0521$, $F^- = 0.0128$)，驗證了群體品味訊號的決策貢獻。
  \item 然而，KG 語意路徑的 $F^- = 0.3134$ 異常偏高，意味著「僅保留 KG 語意路徑」會導致約 31\% 的預測機率損失。我們推測，經過三層非線性訊息傳播後，知識圖譜的結構與語意資訊已被高度壓縮至節點的分散式嵌入向量中，使得單一的顯式路徑無法涵蓋模型實際依賴的全部特徵。此發現可能呼應圖資訊瓶頸假說~\cite{wu2020gib}，並顯示基於注意力的內在可解釋性在深層 GNN 中可能存在結構性的侷限。
  \item 解釋路徑的覆蓋率觀察到約 60\%，約四成的推薦無法透過 Top-$K$ 注意力路徑進行解釋，進一步佐證了分散式嵌入在深層模型決策中的主導地位。
\end{itemize}

\section{研究限制}

本研究仍存在以下限制，需於未來工作中加以改善：

\begin{enumerate}[leftmargin=2em]
  \item \textbf{訓練資源與實驗規模：}受限於深層模型所需的運算時間與記憶體成本，KGAT 目前尚未如 Baseline 模型般進行多次獨立重複運行，且完整指標評估主要集中於偶數 Epoch。此一設定使本文較難充分評估模型效能的統計穩定性，也限制了對長期訓練行為的完整觀察。
  \item \textbf{實驗比較口徑未完全對稱：}Baseline 模型採兩次獨立運行後取最佳 Epoch 平均，而 KGAT 則是在既定資源條件下，以一次訓練過程中的最佳觀察值進行報告。雖然本文已明確說明其原因，但此差異仍意味著 KGAT 與 Baseline 之間的比較應視為在特定條件下的實證觀察，而非已完全排除訓練成本與隨機性的最終結論。
  \item \textbf{單一資料集驗證：}本研究僅在 Food.com 食譜推薦資料集上進行驗證。食譜領域知識圖譜的特殊結構（如高頻樞紐食材節點）可能放大了淺層「資訊汙染」效應。因此，本文所觀察到的深度效應與知識圖譜整合效果，是否能推廣至其他領域資料集，仍有待交叉驗證。
  \item \textbf{評估協定的侷限：}本文採用 Leave-one-out 評估，並為每筆正樣本隨機抽取 100 個負樣本形成候選清單。此作法有助於控制評估成本與比較條件，但也使本文所報告之 HR 與 NDCG 數值主要適合用於相同協定下的相對比較，不宜直接與 full-ranking 或不同負樣本抽樣設定的研究進行絕對值對照。
  \item \textbf{可解釋性分析之樣本有限：}雖然可解釋性評估初始抽取了 100 位使用者，但實際成功萃取出解釋路徑者為 60 位，因此 Fidelity 統計係建立於有效樣本子集合之上。這代表本文所觀察到的忠實度分布，主要反映「可成功萃取路徑之樣本」的特性，而不必然能代表所有推薦結果。
  \item \textbf{Fidelity 指標的固有限制：}本研究採用的傳統 Fidelity 評估透過直接移除或保留邊進行重預測，可能引發資料分佈偏移 (Distribution Shifts)，使評估結果存在一定程度的失真~\cite{agarwal2024fidelity}。因此，本文對解釋路徑充分性與必要性的判讀，仍需結合此方法限制一併理解。
\end{enumerate}

\section{未來研究方向}

基於本研究的發現與上述限制，以下提出四個值得深入探索的未來研究方向：

\begin{enumerate}[leftmargin=2em]
  \item \textbf{擴展訓練規模與重複驗證：}最優先的後續工作為增加 KGAT 模型的訓練 Epoch 數，並進行至少兩次以上的獨立重複運行，報告各指標的平均值與標準差，以確認效能的統計穩定性。同時，引入 DropEdge~\cite{rong2020dropedge}、Message Dropout 或 Early Stopping 等正則化策略，緩解深層模型的訓練不穩定與過擬合風險，進一步挖掘深層架構的潛力。

  \item \textbf{探究 KG 語意路徑充分性偏低的深層機制：}本研究發現 KG 語意路徑的 $F^-$ 顯著高於其他路徑類型。未來研究可從圖資訊瓶頸 (Graph Information Bottleneck)~\cite{wu2020gib} 的角度，分析模型在深層傳播後是否已將 KG 結構資訊壓縮進了分散式嵌入中，從而使得僅保留顯式 KG 路徑不足以支撐預測。此外，可嘗試引入精煉反事實指標 (Refined Counterfactual Metrics) 以緩解傳統 Fidelity 計算中的分佈偏移問題~\cite{agarwal2024fidelity}。

  \item \textbf{跨資料集與跨領域驗證：}為增強研究結論的泛化性，未來應在其他領域的知識圖譜推薦資料集（如 Amazon-Book、Yelp2018 等）上進行交叉驗證，檢驗深度效應、淺層 KG 資訊汙染與 KG 語意路徑充分性危機等結論是否具備跨領域的普適性。

  \item \textbf{消融實驗之系統性擴展：}目前的消融實驗（w/o Attention、w/o KG）僅在 $L{=}1$ 層級進行模組拆解。未來可將消融拓展至 $L{=}2$、$L{=}3$ 等深層配置，觀察注意力機制與知識圖譜在不同深度下的邊際貢獻是否存在質變，以建構更完整的元件貢獻圖像。
\end{enumerate}
